\scanpage{53}
\mspage{22}
$\Delta_n=\{0,\ldots,n\}$ dont l'image soit $A\cup\{n\}$, et l'hom simplicial correspondant
\diagram{\begin{tikzcd}[column sep=small]
\Sigma^{<}(n)\arrow[r,"\mathrm{imm.}"]\arrow[d,phantom,"{\rotatebox{90}{$\in$}}" description]&\Sigma^{<}(d)\arrow[d,phantom,"{\rotatebox{90}{$\in$}}" description]&\deleted{\unsure{(\longrightarrow}}&\big(\arrow[r,dashed,"\mathrm{imm.}"]&\Sigma^{<}(0)\big)\arrow[d,phantom,"{\rotatebox{90}{$\in$}}" description]\\
\widetilde x&z\ \deleted{\unsure{z}}&&&x
\end{tikzcd}}
\diagram{\begin{tikzcd}
\deleted{x}\arrow[r,phantom,"\deleted{\longleftarrow}" description]&\deleted{\unsure{\Sigma(0)}}
\end{tikzcd}}
est l'image de $\widetilde x$ par celui-ci. De plus, pour $z$ dans \deleted{\unsure{ainsi}} $\Sigma^{<}(d)$ au dessus de $x$, \deleted{\unsure{$z$}} le germe de sous-espace de $\Sigma^{<}(0)$ en $x$, image du germe de $\Sigma^{<}(d)$ en $z$, est juste l'intersection des germes $B_\alpha$ correspondants aux $z_\alpha\in A$.

On s'intéresse ici aux morphismes entre les $\Sigma^{<}(n)$ correspondant aux applications strictement croissantes d'ens. tot. ordonnés non vides $d\hookrightarrow d'$ tel que $d$ soit cofinal dans $d'$. La catégorie de ces objets et morphismes équivaut à celle des ens tot. ordonnés $d$ (év\textsuperscript{t} vides), par le foncteur qui à un tel $d$ associe $d\coprod\{\omega\}$, où $\omega$ est un ``dernier élément'', et à tout $d'\hookrightarrow d$ str. croiss le prolongement par $\omega\mapsto\omega$. On a donc \deleted{\unsure{$\Sigma$}} un foncteur
\[
\Gamma:d-\Sigma^{<}(d\coprod\{\omega\})
\]
sur les ens tot. ordonnés év\textsuperscript{t} vides, dont la donnée équivaut \inserted{\unsure{qui est}} à celle des $\Sigma^{<}(n)$ avec les appl. simpliciales ``cofinales'' entre eux. Et la donnée d'un \deleted{\unsure{F}} tel foncteur $\Gamma$ revient à la donnée de $\Gamma(\varnothing)$ (ici $\Sigma^{\unsure{*|<}}(0)$) et
\mspage{52}
